% GENERATED FILE. Edit canonical lessons, then run npm run generate:books.
% canonical-source-hash: fnv1a64:d490c3ec2a6dbd5e
% canonical-lessons: MR-C94-shala, MR-C94-nakasha, MR-C94-dhada, MR-C94-prashna, MR-C94-uttar

\chapter{School, Map, Lesson, Question, Answer}
\label{ch:mr-school-map-lesson-question-answer}
\hlchaptermodality{94}

\begin{chapteropening}
\textbf{By the end of this chapter:} \emph{I can say a school, a map, a lesson, a question and an answer.}

Complete the last lesson of chapter 94: I can say a school, a map, a lesson, a question and an answer.
\end{chapteropening}

\section[\texorpdfstring{\mr{शाळा} (shāḷā)}{शाळा (shāḷā)}]{\mr{शाळा} (shāḷā) — a school}

\label{lesson:MR-C94-shala}



\begin{quote}
Before the new one: say the Marathi for an exercise book, then the Marathi for a pen.
\end{quote}

\subsection*{You'll want to know: \mr{शाळा}}
\textbf{\mr{शाळा}} — \emph{shāḷā} — \textquotedblleft{}a school\textquotedblright{}.

\begin{cousinweb}
\textbf{\mr{शाळा}} (\emph{shāḷā}) is Sanskrit \emph{śālā}, \textquotedblleft{}shed, house\textquotedblright{}, lightly reshaped with Marathi \textbf{\mr{ळ}}. CDIAL lists the inherited Marathi form as \emph{sāḷ}, \textquotedblleft{}workshop, school\textquotedblright{} (CDIAL 12414).
\end{cousinweb}

\begin{grammarlens}[title={What you've built}]
One more word for this chapter.
\end{grammarlens}

\subsection*{Your turn}
\begin{itemize}
\raggedright
  \item \emph{Say it:} \emph{shāḷā}
  \item \emph{Say it:} \emph{shāḷā}, once more
  \item \emph{Say it:} say \emph{shāḷā}
  \item \emph{Recall:} say the Marathi for an exercise book, then the Marathi for a pen, then say \emph{shāḷā} again
\end{itemize}

\begin{tcolorbox}[breakable,colback=teal!4,colframe=teal!35!black,title={Before you move on}]
What does \mr{शाळा} mean? (\textquotedblleft{}A school\textquotedblright{}.) Say it once more.
\end{tcolorbox}
\section[\texorpdfstring{\mr{नकाशा} (nakāshā)}{नकाशा (nakāshā)}]{\mr{नकाशा} (nakāshā) — a map}

\label{lesson:MR-C94-nakasha}



\begin{quote}
Before the new one: say the Marathi for a pen, then the Marathi for a school.
\end{quote}

\subsection*{You'll want to know: \mr{नकाशा}}
\textbf{\mr{नकाशा}} — \emph{nakāshā} — \textquotedblleft{}a map\textquotedblright{}.

\begin{cousinweb}
\textbf{\mr{नकाशा}} (\emph{nakāshā}) is a loan from Arabic \emph{naqsh}, \textquotedblleft{}a painting, an engraving\textquotedblright{}, which reached Marathi through Persian. Its Arabic root is \emph{n-q-sh} (Lane). Persian \emph{naqsha} came to mean a plan or map.
\end{cousinweb}

\begin{grammarlens}[title={What you've built}]
One more word for this chapter.
\end{grammarlens}

\subsection*{Your turn}
\begin{itemize}
\raggedright
  \item \emph{Say it:} \emph{nakāshā}
  \item \emph{Say it:} \emph{nakāshā}, once more
  \item \emph{Say it:} say \emph{nakāshā}
  \item \emph{Recall:} say the Marathi for a pen, then the Marathi for a school, then say \emph{nakāshā} again
\end{itemize}

\begin{tcolorbox}[breakable,colback=teal!4,colframe=teal!35!black,title={Before you move on}]
What does \mr{नकाशा} mean? (\textquotedblleft{}A map\textquotedblright{}.) Say it once more.
\end{tcolorbox}
\section[\texorpdfstring{\mr{धडा} (dhaḍā)}{धडा (dhaḍā)}]{\mr{धडा} (dhaḍā) — a lesson}

\label{lesson:MR-C94-dhada}



\begin{quote}
Before the new one: say the Marathi for a school, then the Marathi for a map.
\end{quote}

\subsection*{You'll want to know: \mr{धडा}}
\textbf{\mr{धडा}} — \emph{dhaḍā} — \textquotedblleft{}a lesson\textquotedblright{}.

\begin{cousinweb}
The origin of \textbf{\mr{धडा}} (\emph{dhaḍā}) is not settled. Neither CDIAL nor DEDR has an entry for it in the sense \textquotedblleft{}lesson\textquotedblright{}.
\end{cousinweb}

\begin{grammarlens}[title={What you've built}]
One more word for this chapter.
\end{grammarlens}

\subsection*{Your turn}
\begin{itemize}
\raggedright
  \item \emph{Say it:} \emph{dhaḍā}
  \item \emph{Say it:} \emph{dhaḍā}, once more
  \item \emph{Say it:} say \emph{dhaḍā}
  \item \emph{Recall:} say the Marathi for a school, then the Marathi for a map, then say \emph{dhaḍā} again
\end{itemize}

\begin{tcolorbox}[breakable,colback=teal!4,colframe=teal!35!black,title={Before you move on}]
What does \mr{धडा} mean? (\textquotedblleft{}A lesson\textquotedblright{}.) Say it once more.
\end{tcolorbox}
\section[\texorpdfstring{\mr{प्रश्न} (prashna)}{प्रश्न (prashna)}]{\mr{प्रश्न} (prashna) — a question}

\label{lesson:MR-C94-prashna}



\begin{quote}
Before the new one: say the Marathi for a map, then the Marathi for a lesson.
\end{quote}

\subsection*{You'll want to know: \mr{प्रश्न}}
\textbf{\mr{प्रश्न}} — \emph{prashna} — \textquotedblleft{}a question\textquotedblright{}.

\begin{cousinweb}
\textbf{\mr{प्रश्न}} (\emph{prashna}) is Sanskrit \emph{praśna}, \textquotedblleft{}question\textquotedblright{}, taken over whole as a learned loan, a tatsama (CDIAL 8818). It is built on \emph{prach}, \textquotedblleft{}to ask\textquotedblright{}.
\end{cousinweb}

\begin{grammarlens}[title={What you've built}]
One more word for this chapter.
\end{grammarlens}

\subsection*{Your turn}
\begin{itemize}
\raggedright
  \item \emph{Say it:} \emph{prashna}
  \item \emph{Say it:} \emph{prashna}, once more
  \item \emph{Say it:} say \emph{prashna}
  \item \emph{Recall:} say the Marathi for a map, then the Marathi for a lesson, then say \emph{prashna} again
\end{itemize}

\begin{tcolorbox}[breakable,colback=teal!4,colframe=teal!35!black,title={Before you move on}]
What does \mr{प्रश्न} mean? (\textquotedblleft{}A question\textquotedblright{}.) Say it once more.
\end{tcolorbox}
\section[\texorpdfstring{\mr{उत्तर} (uttar)}{उत्तर (uttar)}]{\mr{उत्तर} (uttar) — an answer}

\label{lesson:MR-C94-uttar}



\begin{quote}
Before the new one: say the Marathi for a lesson, then the Marathi for a question.
\end{quote}

\subsection*{You'll want to know: \mr{उत्तर}}
\textbf{\mr{उत्तर}} — \emph{uttar} — \textquotedblleft{}an answer\textquotedblright{}.

\begin{cousinweb}
\textbf{\mr{उत्तर}} (\emph{uttar}) is Sanskrit \emph{uttara}, \textquotedblleft{}upper, following\textquotedblright{}, a tatsama; as a noun it already meant \textquotedblleft{}answer\textquotedblright{} in the Ramayana (CDIAL 1767).
\end{cousinweb}

\begin{grammarlens}[title={What you've built}]
One more word for this chapter.
\end{grammarlens}

\subsection*{Your turn}
\begin{itemize}
\raggedright
  \item \emph{Say it:} \emph{uttar}
  \item \emph{Say it:} \emph{uttar}, once more
  \item \emph{Say it:} say \emph{uttar}
  \item \emph{Recall:} say the Marathi for a lesson, then the Marathi for a question, then say \emph{uttar} again
\end{itemize}

\begin{tcolorbox}[breakable,colback=teal!4,colframe=teal!35!black,title={Before you move on}]
What does \mr{उत्तर} mean? (\textquotedblleft{}An answer\textquotedblright{}.) Say it once more.
\end{tcolorbox}
